\documentclass[11pt]{article}
\usepackage{geometry}                % See geometry.pdf to learn the layout options. There are lots.
\geometry{letterpaper}                   % ... or a4paper or a5paper or ... 
%\geometry{landscape}                % Activate for for rotated page geometry
%\usepackage[parfill]{parskip}    % Activate to begin paragraphs with an empty line rather than an indent
\usepackage{graphicx}
\usepackage{amssymb}
\usepackage{epstopdf}
\DeclareGraphicsRule{.tif}{png}{.png}{`convert #1 `dirname #1`/`basename #1 .tif`.png}

\usepackage{algpseudocode}
\usepackage{stmaryrd} % short arrows
% page formatting
\usepackage{fancyhdr}
\pagestyle{fancy}
\cfoot{\thepage}
\rfoot{Copyright \copyright\ Peter Corke 2018}
\lfoot{\texttt{[petercorke.com]}}
\lhead{}

% paragraph formatting
\setlength{\parindent}{0cm}
\setlength{\parskip}{0.5\baselineskip}

% maths notation
\usepackage{amssymb}
\input{rvc-notation}
\usepackage{mathtools}



\title{Manipulator Jacobian based on the elementary transform sequence}
\author{Peter Corke}
\date{January 2015 (\small updated December 2018)}                                           % Activate to display a given date or no date

\begin{document}
\maketitle
%\section{}
%\subsection{}

\section{Introduction}
The Denavit and Hartenberg notation for describing a serial-link mechanism
geometry is a fundamental tool of the roboticist.
Once a manipulator's kinematics is parameterized in this form a large body
of standard algorithms and code implementations for
kinematics, dynamics, motion planning and simulation are available\cite{Corke17a}.

The fundamentals of serial-link robot kinematics and the 
Denavit-Hartenberg\cite{Hartenberg55} notation are well covered in
standard texts\cite{Spong06,Paul81a}.  Each link is represented by
two parameters: the
{\em link length}, $a_i$, and {\em link twist}, $\alpha_i$, which define 
the relative location of the two attached joint axes in space.
The classical method of assigning Denavit-Hartenberg parameters is to 
systematically assign a coordinate frame to each link, but there are
significant constraints on each frame, and also the coordinate system
of the base and the end-effector.
An interesting complication with the Denavit-Hartenberg notation is the
zero-angle configuration, that is, the pose when all joint angles are zero.
For the Puma robot this is a non-obvious `L-shaped' pose with the upper
arm horizontal and the lower arm vertically upward.

Corke\cite{Corke07a} introduced an alternative representation of the kinematics for a serial-link manipulator: the elementary transform sequence (ETS).
This intuitive but systematic approach allows the forward kinematics to be written down by inspection
using a simple ``walk through'' procedure.
The result is a string of elementary translations and rotations, from
the user defined base coordinate to the end-effector.
For example, the ETS for a Puma 560 robot is 
\begin{align*}
T_z(L_1) R_z(q_1) R_y(q_2) T_y(L_2) T_z(L_3) R_y(q_3) T_x(L_6) T_y(L_4) 
T_z(L_5) R_z(q_4) R_y(q_5) R_z(q_6) T_z(L_7) & \label{eq:puma}
\end{align*}
This readable string describes motion along the joint and link chain in terms of translations and rotations
about the canonic axes.  The arguments of the transformations are constants or generalized joint coordinates.
Unlike the Denavit-Hartenberg approach a joint can involve a rotation or translation about \textit{any} axis -- note in this
example the term $R_y(q_2)$ is a rotation about the y-axis.

The paper\cite{Corke07a} goes on to describe an algebraic procedure that can be automatically applied to an ETS
sequence to factorize it into the standard or modified Denavit-Hartenberg
form.  Joint angle offsets, if required by the chosen zero-angle configuration 
chosen, are generated automatically, as are base and tool transformations.
This algorithm is provided in the Robotics Toolbox for \MATLAB\ by the function \var{DHFactor} which is implemented in Java.

Denavit-Hartenberg parameters are still required to determine the manipulator Jacobian and the inverse dynamics.
This article extends \cite{Corke07a} by describing an algorithm to determine the manipulator Jacobian matrix directly from the ETS
representation.

\section{ETS-based forward kinematics}
We will first recap on the forward kinematics, the pose of the robot's end-effector given a set of joint coordinates.

With ETS represent the forward kinematics of a serial-link  robot manipulator by a sequence of $M$ elementary transformations  $\mat{E}_i \in SE(3)$
\begin{equation}
\mat{T}_E = \mat{E}_1 \mat{E}_2 \mat{E}_3 \cdots \mat{E}_M \label{eq:fkine}
\end{equation}
Each factor is either a translation along, or a rotation about, one of either the x-, y- or z-axis.  That is, it is one of six elementary transformations
\[
\mat{E}_i = \left\{ \begin{tabular}{l}
$\mat{T}_x(\eta_i)$ \\
$\mat{T}_y(\eta_i)$ \\
$\mat{T}_z(\eta_i)$ \\
$\mat{R}_x(\eta_i)$ \\
$\mat{R}_y(\eta_i)$ \\
$\mat{R}_z(\eta_i)$ \end{tabular}
\right.
\]
where the parameter is
either a constant (a translational offset or a rotation) or a generalized joint coordinate
\[
\eta_i = \left\{ \begin{tabular}{l}
$c_i$ \\
$q_j$ \end{tabular}
\right.  .
\]
The forward kinematics is simply the cumulative product of the elementary transforms with the joint coordinates plugged in as
appropriate. There is no need to factorize it first into Denavit-Hartenberg parameters.

Using \MATLAB\ we can take a conventional toolbox Denavit-Hartenberg-based model and convert it to ETS format\footnote{This string is different to the one shown above because the link coordinate frames have been assigned using Denavit-Hartenberg constraints.}
\begin{Code}
>> mdl_puma560
>> s = p560.trchain
s =
    'Rz(q1)Rx(90)Rz(q2)Tx(0.4318)Rz(q3)Tz(0.15005)Tx(0.0203)Rx(-90)Rz(q4)
    Tz(0.4318)Rx(90)Rz(q5)Rx(-90)Rz(q6)'
\end{Code}
where translations are in units of metres and angles are in degrees.
We can convert this human-readable string into a sequence of ETS object
\begin{Code}
>> Te = ETS(s)
Te = 
Rz(q1)Rx(90)Rz(q2)Tx(0.4318)Rz(q3)Tz(0.15005)Tx(0.0203)Rx(-90)Rz(q4)Tz(0.4318)
  Rx(90)Rz(q5)Rx(-90)Rz(q6)

>> about Te
Te [ETS] : 1x14 (112 bytes)
\end{Code}
which is a 14-element row vector of \texttt{ETS} objects.  We can look at individual elements in this sequence, for example
\begin{Code}
>> Te(3)
ans = 
Rz(q2)
>> Te(1:3)
ans = 
Rz(q1)Rx(90)Rz(q2)
\end{Code}
and the objects support a number of methods to determine the parameters, joint angle if any and so on.

We can \textit{evaluate} this string for a particular set of joint coordinates
\begin{Code}
>> Te.eval(qz)
ans =
    1.0000         0         0    0.4521
         0    1.0000         0   -0.1500
         0         0    1.0000    0.4318
         0         0         0    1.0000
\end{Code}
which agrees with the classical Denavit-Hartenberg-based forward kinematics
\begin{Code}
>> p560.fkine(qz)
ans = 
         1         0         0    0.4521
         0         1         0     -0.15
         0         0         1    0.4318
         0         0         0         1
\end{Code}
    
\section{ETS-based manipulator Jacobian}
The time derivative of the end-effector pose is 
\begin{equation*}
\dot{\mat{T}_E} = \frac{d \mat{T}_E}{dt} \in \mathbb{R}^{4 \times 4}
\end{equation*}


We are interested in the change of the end-effector pose due to a change in the $j$'th joint so we introduce the joint velocity
\begin{equation}
\frac{d \mat{T}_E}{dt} = \frac{d \mat{T}_E}{d q_j} \frac{d q_j}{dt} = \frac{d \mat{T}_E}{d q_j} \dot{q}_j\label{eq:dTdt}
\end{equation}
The  matrix derivative on the left-hand-side of \eq{eq:dTdt} has 16 elements  -- it is a non-compact representation of the end-effector spatial velocity which is commonly represented by
$(\vec{v},\, \vec{\omega}) \in \mathbb{R}^6$ where $\vec{v}$ is the end-effector translational velocity and $\vec{\omega}$ is the end-effector rotational velocity.
If we write the end-effector pose as
\[
\mat{T}_E = \begin{pmatrix} \mat{R} &\vec{t} \in \SE{3} \\
0\, 0\, 0 & 1 \end{pmatrix}
\]
where $\mat{R} \in \SO{3}$ is an orthonormal rotation matrix representing the orientation of the end-effector and $\vec{t} \in \mathbb{R}^3$ is the translation
of the end-effector then the time derivative  is
\begin{equation}
\dot{\mat{T}} = \begin{pmatrix} \dot{\mat{R}} & \dot{\vec{t}} \\ \label{eq:dTdt2}
0\, 0\, 0 & 0 \end{pmatrix}
\end{equation}


The translational velocity is simply
\[
\vec{v} = \dot{\vec{t}}
\]
Angular velocity is related to the rate of change of a rotation matrix by the well known identity\cite{Corke17a}
\[
 \dot{\mat{R}} = \skx{\vec{\omega}} \mat{R}
 \]
 where $\skx{\cdot} : \mathbb{R}^3 \mapsto \so{3}$ is a skew-symmetric matrix and we can express the angular velocity as
\[
\vec{\omega} = \iskx{ \dot{\mat{R}} \mat{R}^T }
 \]
where $\iskx{\cdot} :  \so{3} \mapsto \mathbb{R}^3 $ is the inverse operator to $\skx{\cdot}$.
Now we can write
 \[
\begin{pmatrix} \vec{v} \\ \vec{\omega} \end{pmatrix} = \begin{pmatrix} \dot{\vec{t}} \\ \iskx{ \dot{\mat{R}} \mat{R}^T} \end{pmatrix} = \vee(\dot{\mat{T}_E}) =  \vee( \frac{d{\mat{T}_E}}{d q_j}) \dot{q}_j
 \]
 where $\isk{\cdot} :  \se{3} \mapsto \mathbb{R}^6 $.


\begin{table}
\centering
\begin{tabular}{|l|l|} \hline
$\frac{d \mat{T}_x(q_j)}{dq_j} = \begin{pmatrix} 0 & 0 & 0 & 1 \\ 0 & 0 & 0 & 0 \\ 0 & 0 & 0 & 0 \\ 0 & 0 & 0 & 0\end{pmatrix} $ & $\frac{d \mat{R}_x(q_j)}{dq_j} =\begin{pmatrix} \skx{1,0,0} \mat{R}_x(q_j) & \begin{matrix} 0 \\ 0 \\ 0 \end{matrix} \\ 0\, 0\, 0 & 0 \end{pmatrix}  $\\
$\frac{d \mat{T}_y(q_j)}{dq_j} =\begin{pmatrix} 0 & 0 & 0 & 0 \\ 0 & 0 & 0 & 1 \\ 0 & 0 & 0 & 0 \\ 0 & 0 & 0 & 0\end{pmatrix} $ & $\frac{d \mat{R}_y(q_j)}{dq_j} =\begin{pmatrix} \skx{0,1,0} \mat{R}_y(q_j) & \begin{matrix} 0 \\ 0 \\ 0 \end{matrix} \\ 0\, 0\, 0 & 0 \end{pmatrix}  $\\
$\frac{d \mat{T}_z(q_j)}{dq_j} =\begin{pmatrix} 0 & 0 & 0 & 0 \\ 0 & 0 & 0 & 0 \\ 0 & 0 & 0 & 1 \\ 0 & 0 & 0 & 0\end{pmatrix} $ & 
$\frac{d \mat{R}_z(q_j)}{dq_j} =\begin{pmatrix} \skx{0,0,1} \mat{R}_z(q_j) & \begin{matrix} 0 \\ 0 \\ 0 \end{matrix} \\ 0\, 0\, 0 & 0 \end{pmatrix}  $\\ \hline
\end{tabular}
\caption{Derivative of primitive transforms with respect to joint coordinate $q_j$.}
\label{tab:derivs}
\end{table}


To determine how motion of the $j$'th joint affects the end-effector pose we can apply the chain rule to \eq{eq:fkine}
\begin{align*}
\frac{d \mat{T}}{dq_j} &= \frac{d \mat{E}_1}{dq_j} \mat{E}_2 \mat{E}_3 \cdots \mat{E}_M \\
&+ \mat{E}_1 \frac{d \mat{E}_2}{dq_j} \mat{E}_3 \cdots \mat{E}_M \\
&+ \mat{E}_1 \mat{E}_2 \frac{d \mat{E}_3}{dq_j}  \cdots \mat{E}_M \\
& \hspace{0.5em} \vdots \\
& + \mat{E}_1 \mat{E}_2 \cdots \frac{d \mat{E}_M}{dq_j} 
\end{align*}
where all the product terms will be zero except for the one where $\mat{E}_i$ is a function of the joint coordinate $q_j$.
If this is the $k$'th term in \eq{eq:fkine}  then we can rewrite \eq{eq:fkine} as
\[
 \mat{T} = \left( \prod_{i=1}^{k-1} \mat{E}_i  \right)  \mat{E}_k    \left( \prod_{i=k+1}^{M} \mat{E}_i  \right)
\]
and the derivative with respect to $q_j$ is
\begin{equation}
\frac{d \mat{T}}{dq_j} = \left( \prod_{i=1}^{k-1} \mat{E}_i  \right)  \frac{d \mat{E}_k}{dq_j}   \left( \prod_{i=k+1}^{M} \mat{E}_i  \right)\label{eq:dtdq}
\end{equation}
The middle derivative term  is easily computed and the value for all possible cases of $\mat{E}_k$ are summarised in Table \ref{tab:derivs}.

The end-effector motion is the sum of the components due to each of the joints
\begin{align}
\begin{pmatrix} \vec{v} \\ \vec{\omega} \end{pmatrix} &=
\vee( \frac{d \mat{T}_E}{d q_1}) \dot{q}_1 + \cdots + \vee( \frac{d \mat{T}_E}{d q_N}) \dot{q}_N 
\end{align}
which can be written in matrix form as
\begin{align}
&= \begin{pmatrix} \vee( \frac{d \mat{T}_E}{d q_1})  \cdots \vee( \frac{d \mat{T}_E}{d q_N})  \end{pmatrix} \begin{pmatrix}\dot{q}_1  \\ \vdots \\ \dot{q}_N \end{pmatrix} \\
&= \mat{J} \dvec{q}
\end{align}
where $\mat{J}$ is the manipulator Jacobian -- the end-effector spatial velocity in the world coordinate frame.
The algorithm is shown succinctly in pseudocode in Figure \ref{fig:algo}.

Continuing the MATLAB example from above
\begin{Code}
>> Te.jacobian(qz)
ans =
    0.1500   -0.4318   -0.4318         0         0         0
    0.4521         0         0         0         0         0
         0    0.4521    0.0203         0         0         0
         0         0         0         0         0         0
         0   -1.0000   -1.0000         0   -1.0000         0
    1.0000         0         0    1.0000         0    1.0000
\end{Code}
which agrees with the classical Denavit-Hartenberg-based algorithm\cite{Paul81c}
\begin{Code}
>> p560.jacob0(qz)
ans =
    0.1500   -0.4318   -0.4318         0         0         0
    0.4521    0.0000    0.0000         0         0         0
         0    0.4521    0.0203         0         0         0
         0         0         0         0         0         0
         0   -1.0000   -1.0000         0   -1.0000         0
    1.0000    0.0000    0.0000    1.0000    0.0000    1.0000
\end{Code}    


\begin{figure}
\begin{center}
\begin{algorithmic}
\For{j = 1 to N}
 \State{$k \gets i$ such that $\eta_i = q_j$}
 \State{$\dmat{T} = \prod_{i=1}^{k-1} \mat{E}_i  \cdot D\left\{ \mat{E}_k \right\} \cdot \prod_{i=k+1}^{M} \mat{E}_i  $}
\State{$\vec{v}$ = $\dmat{T}$[1 $\shortrightarrow$ 3, 4];}
\State{$\dmat{R}$ = $\dmat{T}$[1 $\shortrightarrow$ 3, 1 $\shortrightarrow$ 3];}
\State{$\vec{\omega} = \vex \left( \dmat{R} \mat{R}^T \right)$}
\State{$\mat{J}$ [1 $\shortrightarrow$ 6 \comma j] = $\begin{pmatrix} \vec{v} \comma \vec{\omega} \end{pmatrix}^T$}
\EndFor
\end{algorithmic}
\end{center}
\caption{ETS Jacobian algorithm}\label{fig:algo}
\end{figure}


\section{Conclusion}
Given a robot kinematic model expressed as an elementary transform sequence (ETS), we have shown how to derive the Jacobian
matrix.  This, combined with a standard minimization algorithm, is sufficient to perform numerical inverse kinematics without the nee
for Denavit-Hartenberg parameters.

\bibliographystyle{ieeetr}
\bibliography{strings,robot,kinematics,dynamics,software,publist}
\end{document}  


